\documentclass[10pt,a4paper]{amsart}
\usepackage[margin=19mm]{geometry}
\usepackage{amsmath,amssymb,mathtools}
\usepackage[scr=rsfs,cal=euler]{mathalfa}
\usepackage{xfrac}
\usepackage[unicode,hidelinks,bookmarks=false]{hyperref}
\newcommand{\ie}{i.e.\ }
\newcommand{\cf}{cf.\ }
\newcommand{\resp}{respectively\ }
\newcommand{\Subsection}{\subsection}
\providecommand{\nobreakdash}{\nobreak\hbox{-}}
\setlength{\emergencystretch}{2em}
\multlinegap0pt
\newcommand{\A}{\mathcal{A}}
\theoremstyle{plain}
\newtheorem{theorem}{Theorem}[section]
\newtheorem{lemma}[theorem]{Lemma}
\newtheorem{corollary}[theorem]{Corollary}
\newtheorem{proposition}[theorem]{Proposition}
\theoremstyle{definition}
\newtheorem{definition}[theorem]{Definition}
\newtheorem{example}[theorem]{Example}
\newtheorem{remark}[theorem]{Remark}
\newtheorem{pdef}[theorem]{Proposition-Definition}
\newtheorem{condition}[theorem]{Condition}
\newtheorem{case}{Case:}
\newtheorem{subcase}{Subcase:}
\title[Factorizable Lie conformal bialgebras, quadratic Rota-Baxter Lie conformal algebras and some induced structures]{Factorizable Lie conformal bialgebras, quadratic Rota-Baxter Lie conformal algebras and some induced structures}
\author{Zhongyin Xu, Chengming Bai, Yanyong Hong}
\date{Source: arXiv:2609.28011v1 (CC BY 4.0)}
\begin{document}
\maketitle

\noindent\emph{Source and licence.} The delivered work is the article shown in the title above, version arXiv:2609.28011v1 (CC BY 4.0), distributed under the Creative Commons Attribution 4.0 International licence, as recorded in the arXiv licence metadata for this exact version and shown on the abstract page saved with this item. The full licence notice, which carries the licence URL and the address of that abstract page, is the bundled file \texttt{licenses/arxiv-2609.28011-CC-BY-4.0.txt}.
\par\smallskip

\noindent\textit{Editorial scope.} Counted unit: the article's Theorem (source label flrbl) (source0.tex lines 1039--1052, source label \texttt{flrbl}): ``Let be a factorizable Lie conformal bialgebra with and . Then is a quadratic Rota-Baxter Lie conformal algebra of weight , where the -module homomorphism and are defined respectively by B r - I -1 , I''. It is delivered together with the article's complete original proof (source0.tex lines 1053--1093, one proof environment adjacent to the statement, 41 lines), copied line for line from the article's own text. Required context retained uncounted: I (lines 748--760). Editorial formatting only: an AMS \texttt{amsart} preamble that restates the article's own macro definitions and its own theorem declarations, and the theorem environments the retained lines use; the article's own class file is neither shipped nor input; the retained environments are renumbered sequentially in order of appearance, whereas the article prints them with its own numbering; the article's equation numbering is not reproduced and every cross-reference inside the retained text resolves to the wrapper's numbering. No citation occurs inside any retained line, so the wrapper carries no bibliography; All retained bodies are exact copies of the bundled \texttt{sources/211542/source0.tex}, so no omission receipt applies. No asset-inclusion command occurs inside any retained span and the package ships no image asset.


\section*{Retained context: I (source0.tex lines 748--760)}

\begin{lemma}\label{I}
Let $ (\A,[\cdot_\lambda\cdot])$ be a Lie conformal algebra, where
$\A$ is free as a ${\bf k}[\partial]$-module, and $r\in \A\otimes
\A$. Then the following conclusions hold.
\begin{enumerate}
    \item \label{item1}$I=I^*$. In particular, if $I$ is an isomorphism of Lie conformal algebras, then $(I^{-1})^*=(I^*)^{-1}=I^{-1}$.
    \item \label{item2}The symmetric part of $r$ is $\A$-invariant if and only if
   \begin{equation}\label{eq:rr1}
   	 I\circ {\rm ad}^*(a)_\lambda={\rm ad}(a)_\lambda\circ I,\qquad \forall a\in \A.
   	 \end{equation}
 
\end{enumerate}
\end{lemma}


\section{The counted unit: Theorem (source label flrbl) and its complete original proof}

\begin{theorem}\label{flrbl}
Let $(\A,[\cdot_\lambda \cdot],\Delta_r)$ be a factorizable Lie
conformal bialgebra with $I=r_+-r_-$ and $\theta\in {\bf k}$. Then
$(\A,[\cdot_\lambda\cdot],\mathfrak{B},\omega^I_\lambda(\cdot,\cdot))$
is a quadratic    Rota-Baxter Lie conformal
algebra of weight  $\theta$, where the ${\bf k}[\partial]$-module homomorphism
$\mathfrak{B}:\A\rightarrow \A$ and
$\omega^I_\lambda(\cdot,\cdot)$ are defined respectively by
\begin{align}
    \label{b}       &\mathfrak{B}=\theta r_-\circ I^{-1}, \\
    \label{w}       &\omega^I_\lambda(a,b)=\langle I^{-1}(a),b\rangle_\lambda,  \;\;\;\forall  a,b\in \A.
\end{align}

\end{theorem}

\begin{proof}
Let $a,b\in \A$. Then we have
\begin{eqnarray}
    I([I^{-1}(a)_\lambda I^{-1}(b)]^r)&=&(r_+-r_-)([I^{-1}(a)_\lambda I^{-1}(b)]^r)\label{eq-q1}\\
    &=&[(I+r_-)(I^{-1}(a))_\lambda (I+r_-)(I^{-1}(b)) ]-[r_-(I^{-1}(a))_\lambda r_-(I^{-1}(b))]\nonumber\\
    &=&[a_\lambda b]+[r_-(I^{-1}(a))_\lambda b]+[a_\lambda r_-(I^{-1}(b))].\nonumber
\end{eqnarray}
Thus we obtain
\begin{align*}
    [\mathfrak{B}(a)_\lambda \mathfrak{B}(b)]&=[\theta (r_-\circ I^{-1})(a)_\lambda \theta (r_-\circ I^{-1})(b) ]
    =\theta^2 (r_-\circ I^{-1}\circ I)([I^{-1}(a)_\lambda  I^{-1} (b)]^r)\\
    &=\theta^2 (r_-\circ I^{-1})(I([I^{-1}(a)_\lambda  I^{-1} (b)]^r))\\
    &=\theta^2 (r_-\circ I^{-1})([a_\lambda b]+[r_-(I^{-1}(a))_\lambda b]+[a_\lambda r_-(I^{-1}(b))])\\
    &=\mathfrak{B}(\theta[a_\lambda b]+[\mathfrak{B}(a)_\lambda b]+[a_\lambda
    \mathfrak{B}(b)]).
\end{align*}
Hence $\mathfrak{B}$ is a Rota-Baxter operator of weight $\theta$
on $(\A,[\cdot_\lambda \cdot])$. Obviously,
$\omega^I_\lambda(\cdot,\cdot)$ is non-degenerate. By Lemma
\ref{I}, we have
\begin{align*}
    &\omega_\lambda^I(a,b)=\langle I^{-1}(a),b\rangle_\lambda=\langle b,I^{-1}(a)\rangle_{-\lambda}=\langle I^{-1}(b),a\rangle_{-\lambda}
    =\omega_{-\lambda}^I(b,a),\\
    &\omega_{\lambda}^I([a_\mu b],c)=\langle I^{-1}([a_\mu b]),c\rangle_\lambda=-\langle I^{-1}([b_{-\mu-\partial} a]),c\rangle_\lambda =-\langle ad^* (b)_{\lambda-\mu}(I^{-1}(a)),c\rangle_\lambda\\
    &\qquad\qquad\,\,\,\,=\langle I^{-1}(a),[b_{\lambda-\mu}c]\rangle_\mu=\omega_\mu^I(a,[b_{\lambda-\mu
    }c]).
\end{align*}
Hence $\omega^I_\lambda(\cdot,\cdot) $ is a non-degenerate symmetric invariant
conformal bilinear form on $(\A,[\cdot_\lambda\cdot])$. Note that
$r_-^\ast=-r_+$. Therefore we obtain
\begin{eqnarray*}
    &&\omega_{\lambda}^I(a,\mathfrak{B}(b))+\omega_{\lambda}^I(\mathfrak{B}(a),b)+\theta\omega_{\lambda}^I(a,b)\\
    &=&\langle I^{-1} (a),\theta (r_-\circ I^{-1})(b)\rangle_\lambda +\langle I^{-1}(\theta (r_-\circ I^{-1})a,b\rangle_\lambda +\theta\langle I^{-1}a,b\rangle_\lambda\\
    &=&\theta( \langle I^{-1}(a),(r_-\circ I^{-1})(b)+b\rangle_\lambda+\langle (I^{-1}\circ r_-\circ I^{-1})a,b\rangle_\lambda)\\
    &=&\theta\langle a,(I^{-1}\circ r_-\circ I^{-1}+I^{-1}+I^{-1}\circ (-r_+)\circ I^{-1})b\rangle_\lambda=0.
\end{eqnarray*}
So
$(\A,[\cdot_\lambda\cdot],\mathfrak{B},\omega^I_\lambda(\cdot,\cdot))$
is a quadratic    Rota-Baxter Lie conformal
algebra of weight  $\theta$.
\end{proof}

\end{document}
